\section{Adaptive stopping and complete query counts}
\label{nloracle:section}
The certificates above can be used with ordinary quadratic backtracking.
We state the work bound once for a general calibration $\Phi$; the
specializations then change only this function. Objective values are exact.
Let $D\ge\operatorname{diam}Q>0$, and suppose
\begin{equation}\label{nloracle:model}
 \lambda_{\max}(\nabla^2f(x))\le K(\delta(x)),\qquad
 \delta(x)\le\Phi(G(x)).
\end{equation}
Here $K:[0,\infty)\to[0,\infty)$ is finite, continuous and nondecreasing,
and $\Phi:[0,\infty)\to[0,\infty)$ is continuous,
strictly increasing and unbounded, with $\Phi(0)=0$.
The sharp gap envelope gives $\Phi=\mathcal B$ and $K(u)=H_0+H_1u$.
The known diameter also permits $\Phi=\mathcal B_D$.
These stopping assumptions also allow other regularity conditions.

At requested directional accuracy $\eta>0$, a query returns $h\ge0$
and $d=s-x$, $s\in Q$, such that
\begin{equation}\label{nloracle:query}
 G(x)\le B=h+\eta,\qquad
 -\langle\nabla f(x),d\rangle\ge h-\eta.
\end{equation}
A source certificate $C$ must satisfy $\delta(x)\le C\le\Phi(B)$.
One may use $C=\Phi(B)$ or a sharper valid joint certificate.
Retain $\ell\leftarrow\max\{\ell,f(x)-C\}$ and set $U=f(x)-\ell$.
The initial query creates the first lower bound.

\paragraph{The rule.}
Choose $\eta_0,M>0$ and $0<\theta<1$, and carry $\eta$ between points.
After every query, first stop if $U\le\varepsilon$, where $\varepsilon>0$.
Otherwise require $\eta\le\theta h$. If this fails, halve $\eta$ and
query at the same point, reusing its objective value.
After it passes, put $a=h-\eta$, start with $L=L_{k-1}/2$,
$L_{-1}=M$, and double $L$ until
\begin{equation}\label{eq:btstep}
 \alpha=\min\{1,a/(L\|d\|^2)\},\qquad
 f(x+\alpha d)\le f(x)-\alpha a+\tfrac12L\alpha^2\|d\|^2.
\end{equation}
Accept $x+\alpha d$, cache its value, and retain the accepted $L_k$.
This is the usual backtracking update \cite{pedregosa}; the new
certificates determine when its computed points are accurate enough.

\begin{theorem}[Stopping with carried accuracy]\label{nloracle:main}
Every refinement loop and line search terminates, and the rule stops
with $\delta(x_N)\le U_N\le\varepsilon$.
Let $N$ count accepted updates and $R$ count accuracy halvings, including
initial and terminal source checks. With $\kappa=(1-\theta)/(1+\theta)$,
\begin{equation}\label{nloracle:refinement}
 R\le\left\lceil\log_2\max\left\{1,
 \frac{(1+\theta)\eta_0}{\theta\Phi^{-1}(\varepsilon)}\right\}\right\rceil,
 \qquad T=N+1+R.
\end{equation}
There are exactly $T$ source queries. Each vector query followed by a
full LMO costs one LMO; each score query below costs one full atom scan.

If an initial update is needed, let $U_0>\varepsilon$ be the certificate
after its precision check, and set
\[
 \begin{gathered}
 J=\lceil\log_2(U_0/\varepsilon)\rceil,\quad b_j=U_0/2^j,\quad
 s_j=\kappa\Phi^{-1}(b_j/2),\\
 C_j=\max\{2K(b_j),s_j/D^2\},\quad
 P_K=\left\lceil\log_2\frac{\max\{M,C_0\}}{C_{J-1}}\right\rceil,
 \qquad 0\le j<J.
 \end{gathered}
\]
Then
\begin{equation}\label{nloracle:work}
 N\le T_\varepsilon:=P_K+
 \left\lceil\sum_{j=0}^{J-1}
 \max\left\{\frac{2b_j}{s_j},\frac{4D^2K(b_j)b_j}{s_j^2}\right\}\right\rceil.
\end{equation}
For actual $N\ge1$, all line-search objective calls, including the
initial value and rejected trials, satisfy
\begin{equation}\label{nloracle:functions}
 N_f\le1+2N+
 \left\lceil\log_2\max\{1,2K(U_0)/M\}\right\rceil.
\end{equation}
If $N=0$, there is one objective call and $1+R$ source queries.
The phase quantities are then unnecessary.
\end{theorem}

\begin{proof}
Every retained lower bound is valid, so $0\le\delta\le U\le C\le\Phi(B)$.
At a nonterminal failed precision check,
\[
 \Phi^{-1}(\varepsilon)<\Phi^{-1}(U)\le B
   <(1+\theta)\eta/\theta.
\]
Since the precision is halved and never reset, this proves the bound on
$R$ and finiteness of each check. Each visited point has one closing
query, and each halving has one additional failed query. Hence $T=N+1+R$.
At a passed nonterminal check,
\begin{equation}\label{compact:certified-slope}
 a=h-\eta\ge\kappa(h+\eta)\ge\kappa\Phi^{-1}(U)>0.
\end{equation}

For the line search put $w(t)=f(x+td)-f(x)$ and $K_x=K(\delta(x))$.
While $w\le0$, $w''\le K_x\|d\|^2$ and $w'(0)\le-a<0$.
A first return to zero before $2a/(K_x\|d\|^2)$ contradicts
$w(t)\le-at+K_x\|d\|^2t^2/2<0$. Thus every $L\ge K_x$
passes the test at its proposed step. If $K_x=0$, use the linear
comparison. Every accepted trial satisfies
\begin{equation}\label{eq:acceptedprogress}
 f_k-f_{k+1}\ge\alpha_k a_k/2,\qquad
 L_k\le\max\{L_{k-1}/2,2K(\delta_k)\}.
\end{equation}
The objective decreases and the retained lower bound increases, so $U$
is nonincreasing, also during refinements at a fixed point.

Assign a nonterminal update to $b_j/2<U_k\le b_j$.
The positive phase caps $C_j$ are nonincreasing. If $L_k>C_j$,
the accepted trial must have been the first one: otherwise the preceding
rejected value would be at least $C_j/2\ge K(\delta_k)$.
Such an update therefore halves the inherited curvature.
Before an update in phase $j$, use
\[
 V_k=\max\{0,\log_2(L_{k-1}/C_j)\}.
\]
An update above the cap reduces this potential by one; an update below
the cap leaves zero before a phase change. A jump from phase $j$ to
$j'$ increases it by at most $\log_2(C_j/C_{j'})$. These increments
telescope, so there are at most $P_K$ updates above their caps.

At any remaining update, \eqref{compact:certified-slope} and $\|d\|\le D$
give objective decrease at least
\[
 p_j=\min\{s_j/2,s_j^2/(2C_jD^2)\}=s_j^2/(2C_jD^2).
\]
The total decrease from points in phase $j$ is at most the true error
at its entry, hence at most $b_j$. There are at most $b_j/p_j$ such
updates. Summation proves \eqref{nloracle:work}. The same finite bound
holds for every unfinished prefix, excluding an infinite nonterminal run.

If $q_k$ trials are rejected at update $k$, then
$L_k=2^{q_k-1}L_{k-1}$. Consequently
\[
 \sum_{k<N}(q_k+1)=2N+\log_2(L_{N-1}/M),\qquad
 L_{N-1}\le\max\{M/2^N,2K(U_0)\}.
\]
Adding the initial objective call proves \eqref{nloracle:functions}.
\end{proof}

\paragraph{Away steps.}\label{alg:afw}
For $Q=\operatorname{conv}\mathcal A$ with finitely many atoms, maintain
exact positive weights on the active set $S$.
For AFW, the full LMO returns a minimizing atom in $\mathcal A$, including ties.
Use $B=h_{\rm FW}+\eta$ for stopping and the FW precision check.
Only after that check passes,
scan the active set and compare its approximate away gap with the FW gap.
Select the larger, use FW on ties and singleton sets,
and subtract $\eta$ from the selected gap. The error guarantee must
hold uniformly for both directions, as in the next section.
The away direction $x-v$ has cap $\lambda_v/(1-\lambda_v)$; replace
the cap $1$ in \eqref{eq:btstep} by this value \cite{lacostejaggi}.
A step that reaches the away cap removes an atom. Every other step has
the same progress bound as a FW step. The potential proof counts all
updates above their caps, and the decrease proof counts the remaining
good updates. Thus good updates number at most $T_\varepsilon$.
FW adds at most one atom and a drop removes one, giving
\begin{equation}\label{eq:dropcount}
 N\le2T_\varepsilon+|S_0|-1.
\end{equation}
The query and objective bounds use this actual total $N$; AFW also
performs $N$ active-set scans.

\begin{corollary}[Certified scalar evaluation]\label{nloracle:energy}
Suppose $\delta\le\Phi_0(G)$ and a scalar routine returns
$\Phi_0(B)\le\widehat\Phi_0(B)\le\vartheta\Phi_0(B)$ for known
$\vartheta\ge1$. Use $C=\widehat\Phi_0(B)$, or the smaller of this
value and any other valid source bound. The theorem applies with
$\Phi=\vartheta\Phi_0$. In particular, one may take
$\Phi_0=\mathcal B$ or $\Phi_0=\mathcal B_D$; under $\nabla f(x^*)=0$, one may also take
$\Phi_0=W$. The inverse is needed in the analysis, not in the algorithm.
\end{corollary}
\begin{proof}
Validity gives $\delta\le C$, and construction gives
$U\le C\le\vartheta\Phi_0(B)$. These are the calibration inequalities
used in the proof. The required properties of $\mathcal B$ and $W$
were proved above.
\end{proof}

A norm enclosure $P^+$ permits the exact choice $C=\mathcal J_D(B,P^+)$.
For a computed enclosure, use
$C=\widehat{\mathcal J}_D(B,P^+)\le\vartheta\mathcal J_D(B,P^+)
\le\vartheta\mathcal B_D(B)$ and the effective calibration
$\Phi=\vartheta\mathcal B_D$. The last inequality follows from the
cap characterization of $\mathcal J_D$. Either choice can be combined with other
valid source bounds by taking their minimum. Norm formation and scalar
evaluation are additional work.
If the scalar routine uses $n(B_j)$ comparisons at query $j$, charge
$\sum_{j=1}^{N+1+R}n(B_j)$, including failed refinements and the terminal
query. Comparison counts alone do not bound interval or bit arithmetic.
For a nonincreasing query-cost bound $C_{\rm query}(\eta)$, the total is
at most
\[
 (N+1)C_{\rm query}(\eta_R)+\sum_{r=0}^{R-1}C_{\rm query}(\eta_0 2^{-r}).
\]

\paragraph{Linear gap bounds and the final phases.}
A known $G\ge\gamma\delta$ with $\gamma>0$ gives $\Phi(g)=g/\gamma$.
For $K(u)=H_0+H_1u$ and $\bar\gamma=\kappa\gamma$,
\eqref{nloracle:work} implies
\[
 N\le P_K+\left\lceil
 \frac{4J}{\bar\gamma}+
 \frac{32D^2H_0}{\bar\gamma^2\varepsilon}+
 \frac{16D^2H_1J}{\bar\gamma^2}\right\rceil.
\]
Indeed $s_j=\bar\gamma b_j/2$ and
$\sum_{j<J}b_j^{-1}<2/\varepsilon$.
Quasar-convexity with center $x^*$ supplies this gap bound through
$\langle\nabla f(x),x-x^*\rangle\ge\gamma\delta(x)$.
It need not imply convexity between arbitrary pairs of points.
For the sharp shifted-RSI calibration, the phases with
$b_j/2\le g_c$ have $\mathcal B^{-1}(b_j/2)=b_j/2$ whenever the
positive threshold $g_c$ exists. Their progress bound therefore has
the same form with $\gamma=1$.
For $\Phi=\vartheta\mathcal B$, the corresponding condition is
$b_j/(2\vartheta)\le g_c$, and $s_j=\kappa b_j/(2\vartheta)$.
The same final-phase formulas hold for $\mathcal B_D$, whose ordinary-gap
regime is identical.
